\documentclass[11pt,oneside]{article}

\usepackage[T1]{fontenc}
\usepackage{lmodern}
\usepackage{microtype}
\usepackage[margin=1in,headheight=14pt]{geometry}
\usepackage{amsmath,amssymb,amsthm,mathtools}
\usepackage{booktabs,longtable,tabularx,array}
\usepackage{enumitem}
\usepackage{xcolor}
\usepackage{xurl}
\usepackage{hyperref}
\usepackage{listings}
\usepackage{fancyhdr}
\usepackage{titlesec}
\usepackage{needspace}

\definecolor{proofgreen}{RGB}{15,105,68}
\definecolor{openblue}{RGB}{24,91,145}
\definecolor{computeviolet}{RGB}{98,55,145}
\definecolor{quietgray}{RGB}{86,92,99}
\definecolor{codegray}{RGB}{248,248,248}
\definecolor{rulegray}{RGB}{92,101,111}

\hypersetup{
  colorlinks=true,
  linkcolor=openblue,
  citecolor=proofgreen,
  urlcolor=openblue,
  pdftitle={Centered Cyclotomic Trace Dynamics at the Alaoglu--Erdos Barrier},
  pdfauthor={OpenAI GPT-5.6 Pro, for the ProveIt project},
  pdfsubject={Exact rational superattraction, valuation transport, and fifth-trace cross-gap reductions}
}
\urlstyle{same}

\pagestyle{fancy}
\fancyhf{}
\lhead{Centered cyclotomic trace dynamics}
\rhead{Alaoglu--Erd\H{o}s progress report}
\cfoot{\thepage}
\renewcommand{\headrulewidth}{0.4pt}

\titleformat{\section}{\Large\bfseries}{\thesection}{0.75em}{}
\titleformat{\subsection}{\large\bfseries}{\thesubsection}{0.75em}{}
\titleformat{\subsubsection}{\normalsize\bfseries}{\thesubsubsection}{0.75em}{}

\setlength{\parindent}{0pt}
\setlength{\parskip}{0.55em}
\setlength{\emergencystretch}{3em}
\setlist[itemize]{leftmargin=2em,itemsep=0.25em,topsep=0.35em}
\setlist[enumerate]{leftmargin=2.25em,itemsep=0.25em,topsep=0.35em}

\newtheorem{theorem}{Theorem}[section]
\newtheorem{proposition}[theorem]{Proposition}
\newtheorem{lemma}[theorem]{Lemma}
\newtheorem{corollary}[theorem]{Corollary}
\newtheorem{conjecture}[theorem]{Conjecture}
\newtheorem{criterion}[theorem]{Sufficient criterion}
\theoremstyle{definition}
\newtheorem{definition}[theorem]{Definition}
\newtheorem{experiment}[theorem]{Exact computation}
\newtheorem{target}[theorem]{Research target}
\theoremstyle{remark}
\newtheorem{remark}[theorem]{Remark}
\newtheorem{warning}[theorem]{Caution}

\newcommand{\N}{\mathbb N}
\newcommand{\Z}{\mathbb Z}
\newcommand{\Q}{\mathbb Q}
\newcommand{\R}{\mathbb R}
\newcommand{\ord}{\operatorname{ord}}
\newcommand{\rad}{\operatorname{rad}}
\newcommand{\height}{\operatorname{H}}
\newcommand{\vp}[2]{v_{#1}\!\left(#2\right)}
\newcommand{\Td}{T_d}
\newcommand{\Tp}{T_p}
\newcommand{\thetaAE}{\vartheta}
\newcommand{\statusproved}{\textcolor{proofgreen}{\textbf{[proved here]}}}
\newcommand{\statuscomputed}{\textcolor{computeviolet}{\textbf{[exact computation]}}}
\newcommand{\statusopen}{\textcolor{openblue}{\textbf{[open target]}}}
\newcommand{\statusknown}{\textcolor{quietgray}{\textbf{[background]}}}

\newenvironment{statusbox}[1]
{\begin{center}\begin{minipage}{0.94\textwidth}\raggedright\hrule\vspace{0.42em}
\textbf{Status.} #1\par\vspace{0.32em}}
{\vspace{0.32em}\hrule\end{minipage}\end{center}}

\lstdefinestyle{pythonstyle}{
  language=Python,
  basicstyle=\ttfamily\scriptsize,
  backgroundcolor=\color{codegray},
  frame=single,
  rulecolor=\color{rulegray},
  numbers=left,
  numberstyle=\tiny\color{rulegray},
  numbersep=7pt,
  showstringspaces=false,
  breaklines=true,
  columns=fullflexible,
  keepspaces=true,
  tabsize=4
}

\begin{document}
\pagenumbering{gobble}

\begin{titlepage}
\centering
\vspace*{1.1cm}
{\Huge\bfseries Centered Cyclotomic Trace Dynamics\par}
\vspace{0.2cm}
{\Huge\bfseries at the Alaoglu--Erd\H{o}s Barrier\par}
\vspace{0.7cm}
{\LARGE Exact Rational Superattraction, Prime-Support Transport,\par}
\vspace{0.12cm}
{\LARGE and a Two-Channel Fifth-Trace Reduction\par}
\vspace{1.2cm}
{\large A focused progress addendum to the unified ProveIt report\par}
\vfill
\begin{minipage}{0.90\textwidth}
\small
\textbf{Main outcome.}
Assume hypothetically that a noninteger \(\beta\) has
\(M=2^\beta\) and \(A=3^\beta\) integral.  Every rational approximation
\(p/q\) to \(\beta\) then gives synchronized rational near-units
\[
 z_2=M^q/2^p=2^{q\beta-p},\qquad
 z_3=A^q/3^p=3^{q\beta-p}=z_2^{\log_2 3}.
\]
This addendum studies a new family of arithmetic amplifiers on those near-units,
\[
 T_{2h+1}(z)=\frac1{2h+1}\sum_{j=-h}^{h}z^j.
\]
The cubic trace \(T_3\) is an exact rational superattractor:
\(T_3(z)-1=(z-1)^2/(3z)\).  On reduced fractions it squares the defect numerator,
doubles every nonstructural prime valuation, and introduces no new defect prime.
Its local contact order and rational height degree are both two, so iteration is
\emph{exactly critical}: the extraordinary real contraction is paid for by equally fast
height growth.

The first higher trace, \(T_5\), introduces the positive factor
\(F(a,b)=a^2+3ab+b^2\).  For two reduced rational branches it obeys the new identity
\[
 cdF(a,b)-abF(c,d)=(ad-bc)(ac-bd).
\]
After the structural factor \(5\) is removed, every shared prime of the two
\(T_5\)-defect factors belongs uniquely to one of two channels: the direct cross-gap
\(ad-bc\), or the reciprocal cross-gap \(ac-bd\).  For the Alaoglu--Erd\H{o}s
near-units these become the explicitly synchronized binomial gaps
\[
 \begin{aligned}
  &3^pM^q-2^pA^q,\\
  &(MA)^q-6^p.
 \end{aligned}
\]
Thus the first trace level capable of creating new arithmetic support does not leave an
unstructured cancellation problem; it reduces it to two concrete gap families.

No proof of the conjecture is asserted.  The gain is an exact new mechanism, a complete
height and valuation audit showing why its critical branch stops, and a sharply localized
next target for common-divisor, Kummer, or \(S\)-unit methods.
\end{minipage}
\vfill
{\large August 22, 2026\par}
\end{titlepage}

\pagenumbering{roman}
\begin{abstract}
Let \(\thetaAE=\log_2 3\).  The Alaoglu--Erd\H{o}s conjecture asks whether
\(2^x,3^x\in\Z\) forces \(x\in\Z\).  This report is deliberately narrower than the
supplied unified report: it does not repeat the conditional solution monoid, determinant
hierarchy, Kummer analysis, interpolation theory, carry languages, or finite scans.  Instead
it develops a new centered-cyclotomic trace construction on the synchronized rational
near-units forced by one hypothetical counterexample.

For every odd \(d=2h+1\), the map
\[
 T_d(z)=d^{-1}\sum_{j=-h}^{h}z^j
       =\frac{\sinh(d t/2)}{d\sinh(t/2)}\quad(z=e^t)
\]
has a quadratic fixed-point contact at \(z=1\), a factorization into centered cyclotomic
polynomials, and the Bernoulli--Jordan expansion
\[
 \log T_d(e^t)=\sum_{r\ge1}
 \frac{B_{2r}(d^{2r}-1)}{2r(2r)!}\,t^{2r}.
\]
For a reduced positive fraction \(a/b\), all cancellation in \(T_d(a/b)\) is exactly
\(\gcd(S_d(a,b),d)\), where \(S_d=(a^d-b^d)/(a-b)\); in particular it is bounded by
\(d\).  The exact defect factorization is
\[
 S_d(a,b)-d(ab)^h=(a-b)^2P_d(a,b),
\]
\[
 P_d(a,b)=\sum_{j=1}^{h}(ab)^{h-j}
 \left(\frac{a^j-b^j}{a-b}\right)^2>0.
\]
The rational-map degree is \(d-1\), while the local multiplicity is two.  Consequently every
composition of positive cyclic traces has contact-to-degree ratio at most one, with equality
only for a pure composition of \(T_3\).

The critical map \(T_3\) admits a complete reduced-fraction dynamics.  If
\(z_k=a_k/b_k\), \(u_k=a_k-b_k\), then
\[
 u_{k+1}=u_k^2/3^{\mathbf1_{3\mid u_k}},
\]
and hence, outside \(3\), every valuation doubles exactly and no new prime enters the
defect.  At the same time \(\log H(z_k)=2^k\log H(z_0)+O(2^k)\), while
\(T_3^{\circ k}(e^t)-1=3^{1-2^k}t^{2^k}(1+O_k(t^2))\).  Applied to the synchronized
dyadic and triadic near-units, this produces rational approximants to
\(\thetaAE^{2^k}\), but their relative error remains only \(O_k((q\beta-p)^2)\) and their
height grows on the critical \(2^k\) scale.

For odd primes \(p\ge5\), the structural reduction can also be determined exactly:
\(P_p(a,b)\equiv(a-b)^{p-3}\pmod p\), and if \(p\mid a-b\) then
\(v_p(P_p(a,b))=1\).  Removing that one factor leaves a \(p\)-free new defect factor.
At \(p=5\), this factor is the reciprocal quadratic \(a^2+3ab+b^2\), and the identity above
splits every common prime between two rational branches into direct and reciprocal channels,
with no overlap after structural normalization.  This is the first nontrivial trace level
and the principal new arithmetic reduction of the report.
\end{abstract}

\tableofcontents
\clearpage
\pagenumbering{arabic}

\section{Scope, notation, and the synchronized rational germ}
\label{sec:scope}

\subsection{The conjecture and the narrow purpose of this addendum}

\begin{conjecture}[Alaoglu--Erd\H{o}s]
For every real \(x\),
\[
 2^x\in\Z\quad\hbox{and}\quad3^x\in\Z
 \qquad\Longrightarrow\qquad x\in\Z.
\]
\end{conjecture}

The supplied unified report already gives an extensive conditional structure theory and
pushes determinant, interpolation, valuation, Kummer, compact-orbit, digit, special-function,
and finite-verification methods to explicit barriers.  Repeating those reductions would add
bulk rather than information.  The present report starts from only the minimal hypothetical
data
\begin{equation}
 \beta\notin\Z,\qquad M=2^\beta\in\N,\qquad A=3^\beta\in\N,
 \label{eq:hypothesis}
\end{equation}
and asks whether one can turn rational approximants to \(\beta\) into an arithmetic
contraction stronger than ordinary secants or determinants.

The answer is two-sided.  A centered trace map gives spectacular exact contraction and an
unexpectedly rigid prime-valuation law.  The same exact formulas prove that the best positive
trace is only critical, not supercritical.  The next trace then reveals a new common-divisor
interface which is arithmetically nontrivial but sharply localized.

\begin{statusbox}{All theorem statements labeled \statusproved{} below have complete
self-contained proofs in this document.  Computer tables are independently regenerated by
the exact-integer script in Appendix~\ref{app:code}.  Statements explicitly called targets
are not used as theorems.}
\end{statusbox}

\subsection{A counterexample supplies rational points approaching one}

A nonintegral rational \(\beta=r/s\) in lowest terms cannot satisfy
\(2^\beta\in\Z\): from \((2^\beta)^s=2^r\), unique factorization forces \(s\mid r\).
Thus every hypothetical \(\beta\) in~\eqref{eq:hypothesis} is irrational, without invoking
transcendence theory.

For integers \(q\ge1\) and \(p\), put
\begin{equation}
 \delta_{p,q}=q\beta-p,
 \qquad z_{2;p,q}=\frac{M^q}{2^p},
 \qquad z_{3;p,q}=\frac{A^q}{3^p}.
 \label{eq:synchronized-germ}
\end{equation}

\begin{proposition}[Synchronized rational germ]
\label{prop:germ}
Under~\eqref{eq:hypothesis},
\begin{equation}
 z_{2;p,q}=2^{\delta_{p,q}}\in\Q_{>0},\qquad
 z_{3;p,q}=3^{\delta_{p,q}}\in\Q_{>0},\qquad
 z_{3;p,q}=z_{2;p,q}^{\thetaAE},
 \quad \thetaAE=\frac{\log3}{\log2}.
 \label{eq:germ-relations}
\end{equation}
Along the continued-fraction convergents \(p/q\) to \(\beta\), both rational numbers tend to
one, from alternating sides, and they use the same infinitesimal parameter \(\delta_{p,q}\).
\end{proposition}

\begin{proof}
The first two identities follow from \(M=2^\beta\) and \(A=3^\beta\).  The third is
\((2^\delta)^{\log_2 3}=3^\delta\).  Convergents satisfy
\(q\beta-p\to0\), with alternating sign for an irrational \(\beta\).
\end{proof}

The essential point is arithmetic, not merely analytic: the same small real parameter is
encoded by two rational numbers whose reduced denominators are supported at the structural
primes two and three.  We shall apply rational maps to the pair
\((z_{2;p,q},z_{3;p,q})\) and audit both the real contraction and the exact reduced
numerators.

\subsection{How this differs from the existing secant construction}

The unified report studies the rational symmetric secant
\(\sinh(xh)/\sinh h\), where the unknown exponent \(x\) occupies the dilation slot and the
small parameter is a \(2,3\)-unit.  Here the dilation is an auxiliary \emph{odd integer}
\(d\), while the rational near-unit \(z\) is the dynamical variable.  This reversal has two
consequences absent from the secant formulation:
\begin{enumerate}
\item the map can be iterated as a rational self-map of \(\mathbb P^1\);
\item reduction of every iterate can be followed prime by prime using integer polynomial
identities.
\end{enumerate}
The analytic kernel is related, but the arithmetic problem and the resulting recurrences are
different.

\section{Centered cyclic traces and their cyclotomic resolution}
\label{sec:trace}

\subsection{Definition and elementary real geometry}

\begin{definition}[Centered cyclic trace]
For an odd integer \(d=2h+1\ge3\), define
\begin{equation}
 T_d(z)=\frac1d\sum_{j=-h}^{h}z^j
 =\frac{z^{-h}(z^d-1)}{d(z-1)},
 \label{eq:Td-definition}
\end{equation}
with the removable value \(T_d(1)=1\).
\end{definition}

The terminology comes from averaging the characters of the \(d\)-dimensional cyclic weight
string \(-h,-h+1,\ldots,h\).  It is also a normalized, centered \(q\)-integer.

\begin{proposition}[Reciprocity, positivity, and hyperbolic form]
\label{prop:basic-trace}
For \(z>0\),
\begin{align}
 T_d(z^{-1})&=T_d(z),\label{eq:reciprocity}\
 T_d(z)&\ge1,\quad\hbox{with equality exactly at }z=1,\label{eq:AMGM}\
 T_d(e^t)&=\frac{\sinh(dt/2)}{d\sinh(t/2)}.\label{eq:hyperbolic}
\end{align}
In particular, \(z=1\) is a fixed point of local multiplicity two.
\end{proposition}

\begin{proof}
Reciprocity reverses the summation index.  The product of the \(d\) positive summands
\(z^{-h},\ldots,z^h\) is one, so their arithmetic mean is at least one, with equality
only when every summand is one.  Summing the finite geometric progression after writing
\(z=e^t\) gives~\eqref{eq:hyperbolic}.  The right side is an even analytic function of
\(t\), with nonzero quadratic term by Theorem~\ref{thm:bernoulli-trace} below.
\end{proof}

For \(d=3\) there are two identities worth displaying immediately:
\begin{equation}
 T_3(z)=\frac{z^2+z+1}{3z},
 \qquad
 T_3(z)-1=\frac{(z-1)^2}{3z}.
 \label{eq:T3-exact}
\end{equation}
Thus the fixed point is not merely quadratic in a Taylor approximation: its defect is an
exact square in the rational function field.

\subsection{The Bernoulli expansion}

Let \(B_n\) denote the Bernoulli numbers with \(B_2=1/6\), \(B_4=-1/30\).

\begin{theorem}[Exact logarithmic expansion]
\label{thm:bernoulli-trace}
As a convergent expansion for \(|t|<2\pi/d\), and also as a formal power series,
\begin{equation}
 \boxed{
 \log T_d(e^t)=\sum_{r\ge1}
 \frac{B_{2r}(d^{2r}-1)}{2r(2r)!}\,t^{2r}.}
 \label{eq:Bernoulli-trace}
\end{equation}
In particular,
\begin{equation}
 \log T_d(e^t)=\frac{d^2-1}{24}t^2
 -\frac{d^4-1}{2880}t^4+O_d(t^6).
 \label{eq:trace-first-terms}
\end{equation}
\end{theorem}

\begin{proof}
The classical formal identity
\[
 \log\frac{e^u-1}{u}=\frac u2+
 \sum_{r\ge1}\frac{B_{2r}}{2r(2r)!}u^{2r}
\]
gives
\[
 \log T_d(e^t)=-\frac{d-1}{2}t
 +\log\frac{e^{dt}-1}{dt}-\log\frac{e^t-1}{t}.
\]
The linear terms cancel, and the even coefficients are exactly those in
\eqref{eq:Bernoulli-trace}.
\end{proof}

The cancellation of the linear term is the analytic source of the quadratic defect.  The
next subsection identifies its arithmetic source.

\subsection{Centered cyclotomic factors and Jordan totients}

For \(n>1\), define the centered cyclotomic factor
\begin{equation}
 \mathscr C_n(z)=z^{-\varphi(n)/2}\Phi_n(z).
 \label{eq:centered-cyclotomic}
\end{equation}
When \(n\) is odd, \(\varphi(n)\) is even and this is a rational Laurent polynomial invariant
under \(z\mapsto z^{-1}\).

\begin{theorem}[Cyclotomic resolution of the trace]
\label{thm:cyclotomic-resolution}
For every odd \(d\ge3\),
\begin{equation}
 \boxed{
 T_d(z)=\frac1d\prod_{\substack{n\mid d\\n>1}}\mathscr C_n(z).}
 \label{eq:trace-cyclotomic-factorization}
\end{equation}
Moreover, with the Jordan totient
\[
 J_m(n)=\sum_{e\mid n}\mu(n/e)e^m,
\]
one has
\begin{equation}
 \log\mathscr C_n(e^t)
 =\log\Phi_n(1)+
 \sum_{r\ge1}\frac{B_{2r}J_{2r}(n)}{2r(2r)!}t^{2r}.
 \label{eq:Jordan-cyclotomic-expansion}
\end{equation}
Summing~\eqref{eq:Jordan-cyclotomic-expansion} over \(n\mid d\), \(n>1\), recovers
Theorem~\ref{thm:bernoulli-trace} from the divisor identity
\(\sum_{n\mid d}J_m(n)=d^m\).
\end{theorem}

\begin{proof}
The ordinary factorization
\(z^d-1=(z-1)\prod_{n\mid d,n>1}\Phi_n(z)\), together with
\(\sum_{n\mid d,n>1}\varphi(n)=d-1\), gives
\eqref{eq:trace-cyclotomic-factorization}.  For the expansion, use
\[
 \Phi_n(e^t)=\prod_{e\mid n}(e^{et}-1)^{\mu(n/e)}.
\]
Because \(n>1\), the sum of the M\"obius exponents is zero.  The constant product is
\(\Phi_n(1)\), the centered linear term is canceled by
\(-\varphi(n)t/2\), and the coefficient of \(t^{2r}\) is the M\"obius sum
\(J_{2r}(n)\).  Finally, summing Jordan totients over divisors gives \(d^{2r}-1\) after
removing the divisor one.
\end{proof}

\begin{remark}[Why this is a Witt/necklace-adjacent construction]
The necklace polynomial
\[
 \mathcal N_n(X)=\frac1n\sum_{e\mid n}\mu(e)X^{n/e}
\]
is integer-valued on integers and is a standard coordinate in the necklace/Witt formalism
\cite{MetropolisRota1983}.  But for \(n>1\),
\[
 \mathcal N_n(1)=0,
 \qquad \mathcal N_n'(1)=\sum_{e\mid n}\frac{\mu(e)}e=\frac{\varphi(n)}n\ne0.
\]
Raw necklace coordinates therefore have only first-order contact at one.  Centering by
reciprocity is exactly what removes that first-order term.  The trace maps are not being
presented as a new construction of the Witt ring; rather, the Witt/necklace factorization
explains why the cyclotomic pieces and Jordan totients appear, and why ordinary ghost
coordinates do not supply the desired superattraction.
\end{remark}

\section{Exact rational arithmetic of the trace maps}
\label{sec:rational-arithmetic}

\subsection{All denominator cancellation is structural}

For odd \(d=2h+1\), put
\begin{equation}
 S_d(a,b)=\frac{a^d-b^d}{a-b}
 =a^{d-1}+a^{d-2}b+\cdots+b^{d-1}.
 \label{eq:Sd}
\end{equation}

\begin{theorem}[Exact reduced denominator]
\label{thm:exact-denominator}
Let \(a,b\in\N\) be coprime.  Then
\begin{equation}
 T_d(a/b)=\frac{S_d(a,b)}{d(ab)^h}.
 \label{eq:rational-Td}
\end{equation}
Furthermore
\begin{equation}
 \gcd(S_d(a,b),ab)=1,
 \qquad
 \gcd\bigl(S_d(a,b),d(ab)^h\bigr)=g_d(a,b),
 \label{eq:exact-gd}
\end{equation}
where
\begin{equation}
 g_d(a,b)=\gcd(S_d(a,b),d)\mid d.
 \label{eq:gd-definition}
\end{equation}
Thus reduction of \(T_d(a/b)\) can cancel only a divisor of the fixed integer \(d\); no
prime carried by \(a\) or \(b\) cancels.
\end{theorem}

\begin{proof}
Substitution into~\eqref{eq:Td-definition} gives~\eqref{eq:rational-Td}.  Modulo \(a\), the
sum \(S_d(a,b)\) is \(b^{d-1}\), and modulo \(b\) it is \(a^{d-1}\).  Coprimality of
\(a,b\) therefore gives the first identity in~\eqref{eq:exact-gd}; the second follows
immediately.
\end{proof}

This exact audit is stronger than a generic height estimate.  In particular, any hoped-for
large cancellation in the trace denominator must occur only after two branches are compared;
it cannot be hidden inside one trace value.

\subsection{A positive sum-of-squares defect factor}

\begin{theorem}[Exact defect factorization]
\label{thm:defect-factorization}
For odd \(d=2h+1\ge3\), define
\begin{equation}
 P_d(a,b)=\sum_{j=1}^{h}(ab)^{h-j}
 \left(\frac{a^j-b^j}{a-b}\right)^2.
 \label{eq:Pd-sos}
\end{equation}
Then \(P_d\in\Z[a,b]\) is a symmetric, positive homogeneous polynomial of degree \(d-3\),
and
\begin{equation}
 \boxed{
 S_d(a,b)-d(ab)^h=(a-b)^2P_d(a,b).}
 \label{eq:defect-factorization}
\end{equation}
If \(T_d(a/b)=a'/b'\) is reduced and \(u=a-b\), then
\begin{equation}
 a'-b'=\frac{u^2P_d(a,b)}{g_d(a,b)}.
 \label{eq:reduced-defect-general}
\end{equation}
\end{theorem}

\begin{proof}
From reciprocity,
\begin{align*}
 T_d(z)-1
 &=\frac1d\sum_{j=1}^{h}(z^j+z^{-j}-2)\\
 &=\frac1d\sum_{j=1}^{h}\frac{(z^j-1)^2}{z^j}.
\end{align*}
Substituting \(z=a/b\), extracting \((a-b)^2\), and putting the terms over the common
denominator \((ab)^h\) gives
\[
 T_d(a/b)-1=\frac{(a-b)^2P_d(a,b)}{d(ab)^h}.
\]
Comparison with~\eqref{eq:rational-Td} proves~\eqref{eq:defect-factorization}; division by
the exact reduction factor \(g_d\) proves~\eqref{eq:reduced-defect-general}.  Positivity on
positive \(a,b\) is manifest in~\eqref{eq:Pd-sos}.
\end{proof}

Expanding the sum of squares gives a useful coefficient formula:
\begin{equation}
 P_{2h+1}(a,b)=
 \sum_{r=0}^{2h-2}
 \binom{\min(r,2h-2-r)+2}{2}a^{2h-2-r}b^r.
 \label{eq:Pd-coefficients}
\end{equation}
The first cases are
\begin{align}
 P_3(a,b)&=1,\label{eq:P3}\
 P_5(a,b)&=a^2+3ab+b^2,\label{eq:P5}\
 P_7(a,b)&=a^4+3a^3b+6a^2b^2+3ab^3+b^4.\label{eq:P7}
\end{align}
Thus \(T_3\) is the unique member of the family whose defect is a pure square up to the
fixed structural reduction.

\subsection{Exact height degree}

For a reduced positive rational \(z=a/b\), write
\(H(z)=\max(a,b)\) and \(h(z)=\log H(z)\).

\begin{theorem}[Height growth of one trace]
\label{thm:height-growth}
For coprime positive \(a,b\) and odd \(d\ge3\),
\begin{equation}
 \frac{H(a/b)^{d-1}}d
 \le H\bigl(T_d(a/b)\bigr)
 \le dH(a/b)^{d-1}.
 \label{eq:height-growth}
\end{equation}
Equivalently,
\begin{equation}
 \left|h(T_d(z))-(d-1)h(z)\right|\le\log d.
 \label{eq:log-height-growth}
\end{equation}
\end{theorem}

\begin{proof}
The \(d\) monomials in \(S_d\) have geometric mean \((ab)^h\), hence
\(S_d\ge d(ab)^h\).  The reduced numerator therefore dominates the reduced denominator, so
\(H(T_d(a/b))=S_d/g_d\).  If \(H=\max(a,b)\), then
\(H^{d-1}\le S_d\le dH^{d-1}\), while \(1\le g_d\le d\).
\end{proof}

The analytic contact at one is always two; the arithmetic degree is \(d-1\).  This gives a
sharp global comparison.

\begin{theorem}[Unique criticality of the cubic trace]
\label{thm:criticality}
Let
\[
 \Psi=T_{d_m}\circ\cdots\circ T_{d_1},
 \qquad d_i\ge3\text{ odd}.
\]
Then the local multiplicity of \(\Psi(z)-1\) at \(z=1\) is \(2^m\), while the degree of
\(\Psi\) as a rational map of \(\mathbb P^1\) is
\begin{equation}
 \deg\Psi=\prod_{i=1}^{m}(d_i-1)\ge2^m.
 \label{eq:composition-degree}
\end{equation}
Equality holds exactly when \(d_1=\cdots=d_m=3\).
\end{theorem}

\begin{proof}
Each \(T_d(z)-1\) has a zero of exact order two by
Theorem~\ref{thm:bernoulli-trace}; local multiplicities multiply under composition.  The
reduced rational function \(T_d\) has numerator degree \(d-1\), denominator degree \(h\),
and no common factor, so its map degree is \(d-1\).  Rational-map degrees also multiply.
Since every \(d_i-1\ge2\), the final assertion follows.
\end{proof}

\begin{remark}[The contact/height verdict]
A rational number differing from one is separated from one by at least the reciprocal of its
reduced denominator.  Theorem~\ref{thm:criticality} says that no composition of positive
cyclic traces can make local contact grow faster than rational degree.  The all-cubic
composition reaches equality, so it deserves a complete arithmetic analysis; every use of a
higher positive trace is subcritical before any cross-branch cancellation is considered.
\end{remark}

\section{The exact arithmetic dynamics of \texorpdfstring{$T_3$}{T3}}
\label{sec:T3}

\subsection{Reduced recurrence}

Let \(z_0=a_0/b_0\in\Q_{>0}\) be reduced and define
\(z_{k+1}=T_3(z_k)=a_{k+1}/b_{k+1}\) in lowest terms.  Put
\begin{equation}
 u_k=a_k-b_k,
 \qquad
 g_k=\gcd(a_k^2+a_kb_k+b_k^2,3).
 \label{eq:T3-data}
\end{equation}

\begin{theorem}[Exact reduced cubic-trace dynamics]
\label{thm:T3-recurrence}
For every \(k\ge0\),
\begin{align}
 a_{k+1}&=\frac{a_k^2+a_kb_k+b_k^2}{g_k},
 &b_{k+1}&=\frac{3a_kb_k}{g_k},\label{eq:T3-ab-recurrence}\\
 u_{k+1}&=\frac{u_k^2}{g_k},
 &g_k&=3^{\mathbf1_{3\mid u_k}}.\label{eq:T3-u-recurrence}
\end{align}
In particular, after the first step \(u_k\ge0\), and it is positive unless \(z_0=1\).
\end{theorem}

\begin{proof}
The first line is Theorem~\ref{thm:exact-denominator} at \(d=3\).  Subtracting the two
reduced coordinates and using
\[
 a_k^2+a_kb_k+b_k^2-3a_kb_k=(a_k-b_k)^2
\]
gives the defect recurrence.  Finally,
\[
 a_k^2+a_kb_k+b_k^2=(a_k-b_k)^2+3a_kb_k,
\]
so the numerator is divisible by three exactly when \(u_k\) is.
\end{proof}

\subsection{Closed-form valuation transport}

\begin{theorem}[Exact prime-valuation law]
\label{thm:T3-valuations}
Let \(u_0\ne0\).  For every prime \(\ell\ne3\),
\begin{equation}
 v_\ell(u_k)=2^k v_\ell(u_0).
 \label{eq:non3-valuations}
\end{equation}
At the structural prime three,
\begin{equation}
 v_3(u_k)=
 \begin{cases}
 0,&v_3(u_0)=0,\\
 1+2^k\bigl(v_3(u_0)-1\bigr),&v_3(u_0)\ge1.
 \end{cases}
 \label{eq:3-valuation}
\end{equation}
Equivalently, for \(k\ge1\),
\begin{equation}
 |u_k|=
 \begin{cases}
 |u_0|^{2^k},&3\nmid u_0,\\
 3\left(|u_0|/3\right)^{2^k},&3\mid u_0.
 \end{cases}
 \label{eq:T3-closed-form}
\end{equation}
\end{theorem}

\begin{proof}
For \(\ell\ne3\), equation~\eqref{eq:T3-u-recurrence} doubles the valuation.  If the
initial three-adic valuation is zero, it remains zero and no factor three is removed.  If it
is positive, the recurrence is \(s_{k+1}=2s_k-1\), whose solution is
\(s_k=1+2^k(s_0-1)\).  Recombining all prime valuations gives
\eqref{eq:T3-closed-form}.
\end{proof}

This law is unusually rigid.  Outside three, the support of the numerator defect is exactly
constant and each existing valuation doubles.  At three, the support is again constant and
only one valuation unit is lost at each squaring stage.

\begin{corollary}[No new cross-orbit support]
\label{cor:no-new-support}
Let \(u_k\) and \(v_k\) be the defect numerators of two reduced \(T_3\)-orbits.  Then
\begin{equation}
 \rad_{\ne3}\gcd(u_k,v_k)
 =\rad_{\ne3}\gcd(u_0,v_0),
 \label{eq:radical-invariance}
\end{equation}
and for every \(\ell\ne3\),
\begin{equation}
 v_\ell\bigl(\gcd(u_k,v_k)\bigr)
 =2^k v_\ell\bigl(\gcd(u_0,v_0)\bigr).
 \label{eq:gcd-valuation-doubling}
\end{equation}
The structural prime three has the corresponding explicit law obtained by taking the minimum
in~\eqref{eq:3-valuation}.
\end{corollary}

\begin{proof}
Take the minimum of the two valuation formulas in
Theorem~\ref{thm:T3-valuations}.
\end{proof}

Thus cubic iteration can magnify common arithmetic support already present in the two initial
gaps, but it cannot manufacture a new common prime.  This is the exact reason that the most
attractive trace map cannot by itself create the missing cross-base cancellation.

\subsection{Real superattraction versus height growth}

For \(z>1\), equation~\eqref{eq:T3-exact} gives \(T_3(z)>1\), while
\begin{equation}
 T_3(z)-z=-\frac{(2z+1)(z-1)}{3z}<0.
 \label{eq:T3-monotone}
\end{equation}
Every positive real point other than one is therefore mapped above one and then decreases to
one.

\begin{theorem}[Archimedean basin and arithmetic escape]
\label{thm:basin-height}
For every \(z_0\in\Q_{>0}\setminus\{1\}\),
\begin{equation}
 T_3^{\circ k}(z_0)\longrightarrow1
 \qquad(k\to\infty),
 \label{eq:real-convergence}
\end{equation}
but its rational height tends to infinity doubly exponentially in the iteration count.  More
precisely, if \(H_k=H(z_k)\), then
\begin{equation}
 \frac{H_k^2}{3}\le H_{k+1}\le3H_k^2,
 \label{eq:T3-height-recurrence}
\end{equation}
and
\begin{equation}
 2^k\log H_0-(2^k-1)\log3
 \le\log H_k\le
 2^k\log H_0+(2^k-1)\log3.
 \label{eq:T3-height-iterate}
\end{equation}
The limit
\begin{equation}
 \widehat h_3(z_0)=\lim_{k\to\infty}2^{-k}\log H_k
 \label{eq:canonical-height}
\end{equation}
exists, is positive for \(z_0\ne1\), and satisfies
\begin{equation}
 \left|\widehat h_3(z_0)-\log H_0\right|
 \le \log 3.
 \label{eq:canonical-height-bound}
\end{equation}
\end{theorem}

\begin{proof}
The real convergence follows from~\eqref{eq:T3-monotone} and the positive fixed-point
equation.  The height recurrence is Theorem~\ref{thm:height-growth} at \(d=3\), and its
iteration gives~\eqref{eq:T3-height-iterate}.  The bounded difference
\(|\log H_{k+1}-2\log H_k|\le\log3\) makes \(2^{-k}\log H_k\) Cauchy and proves the
bound on its limit.  For a nontrivial reduced positive fraction, the first image has height at
least seven; applying the lower recurrence from that point gives a positive limiting lower
bound.
\end{proof}

The map is therefore an explicit arithmetic-dynamical paradox only in appearance: it is a
superattractor at the real place, while its global height escapes with exactly the map degree.
The product formula pays for the real contraction at other places and in the growing
numerator/denominator support.

\subsection{Exact error recurrence and critical saturation}

Put \(\varepsilon_k=z_k-1\) after the first iterate, so \(\varepsilon_k>0\).  Then
\begin{equation}
 \varepsilon_{k+1}=\frac{\varepsilon_k^2}{3(1+\varepsilon_k)}.
 \label{eq:epsilon-recurrence}
\end{equation}
If \(0<\varepsilon_k\le1\),
\begin{equation}
 \frac{\varepsilon_k^2}{6}
 \le\varepsilon_{k+1}\le
 \frac{\varepsilon_k^2}{3}.
 \label{eq:epsilon-bounds}
\end{equation}
At the same time, because \(\varepsilon_k=(a_k-b_k)/b_k\) is reduced as a rational
difference,
\begin{equation}
 \varepsilon_k\ge\frac1{b_k}\ge\frac1{H_k}.
 \label{eq:rational-separation}
\end{equation}
The superattraction and the rational-separation floor therefore live on the same \(2^k\)
logarithmic scale.

\begin{proposition}[Leading local iterate]
\label{prop:T3-local-iterate}
For every fixed \(k\ge1\),
\begin{equation}
 T_3^{\circ k}(e^t)-1
 =3^{1-2^k}t^{2^k}\bigl(1+O_k(t^2)\bigr).
 \label{eq:T3-local-iterate}
\end{equation}
More precisely, if the relative quadratic correction is denoted by \(\alpha_k\), then
\begin{equation}
 \alpha_1=\frac1{12},
 \qquad
 \alpha_k=-\frac{2^{k-3}}3\quad(k\ge2).
 \label{eq:alpha-k}
\end{equation}
\end{proposition}

\begin{proof}
The first step is
\[
 T_3(e^t)-1=\frac{t^2}{3}\left(1+\frac{t^2}{12}+O(t^4)\right).
\]
Substitute an expansion \(c_kt^{2^k}(1+\alpha_kt^2+O_k(t^4))\) into
\eqref{eq:epsilon-recurrence}.  The leading constants satisfy
\(c_{k+1}=c_k^2/3\), giving \(c_k=3^{1-2^k}\).  At the transition from \(k=1\) to
\(k=2\), the denominator contributes at relative order \(t^2\), producing
\(\alpha_2=-1/6\); afterward that denominator correction starts at order at least \(t^4\),
so \(\alpha_{k+1}=2\alpha_k\).
\end{proof}

\subsection{Exact numerical illustration}

Table~\ref{tab:T3-orbits} was generated using only integer recurrence
\eqref{eq:T3-ab-recurrence}.  The first orbit begins with unit defect and nearly saturates
\(-\log_{10}|z_k-1|\le\log_{10}H_k\) at every stage.  The second begins with defect three;
Theorem~\ref{thm:T3-valuations} keeps the reduced defect equal to three forever.

\begin{table}[htbp]
\centering
\small
\begin{tabular}{@{}c r r r r@{}}
\toprule
Initial fraction & \(k\) & \(\log_{10}H_k\) &
\(-\log_{10}|z_k-1|\) & \(|u_k|\)\\
\midrule
\(1000001/1000000\) &0&6.000000&6.000000&1\\
&1&12.477122&12.477121&1\\
&2&25.431364&25.431364&1\\
&3&51.339849&51.339849&1\\
&4&103.156819&103.156819&1\\
\addlinespace
\(1000003/1000000\) &0&6.000001&5.522879&3\\
&1&12.477124&12.000002&3\\
&2&25.431369&24.954247&3\\
&3&51.339858&50.862737&3\\
&4&103.156837&102.679716&3\\
\bottomrule
\end{tabular}
\caption{Exact cubic-trace orbits.  The decimal columns are displays of exact integer
fractions, not floating-point evidence for any theorem.}
\label{tab:T3-orbits}
\end{table}

\section{Applying the cubic trace to a hypothetical counterexample}
\label{sec:synchronized-T3}

Fix \(p,q\) and abbreviate the synchronized germ of Proposition~\ref{prop:germ} by
\[
 z_2=2^\delta,\qquad z_3=3^\delta,
 \qquad \delta=q\beta-p.
\]
For \(k\ge1\), put
\begin{equation}
 E_{2,k}=T_3^{\circ k}(z_2)-1,
 \qquad
 E_{3,k}=T_3^{\circ k}(z_3)-1.
 \label{eq:synchronized-errors}
\end{equation}
Both are positive rationals, even when \(\delta<0\), because the first trace is reciprocal
and strictly above one off the fixed point.

\subsection{A tower of rational approximants to powers of \texorpdfstring{$\thetaAE$}{theta}}

\begin{theorem}[Trace-tower approximants]
\label{thm:tower-approximants}
For every fixed \(k\ge1\),
\begin{equation}
 R_{k;p,q}:=\frac{E_{3,k}}{E_{2,k}}\in\Q_{>0},
 \label{eq:Rkpq}
\end{equation}
and, as \(\delta\to0\),
\begin{equation}
 R_{k;p,q}
 =\thetaAE^{2^k}
 \left(1+\alpha_k\bigl((\log3)^2-(\log2)^2\bigr)\delta^2
 +O_k(\delta^4)\right),
 \label{eq:tower-approximation}
\end{equation}
where \(\alpha_k\) is given by~\eqref{eq:alpha-k}.
\end{theorem}

\begin{proof}
Rationality follows because \(T_3\in\Q(z)\).  Apply
Proposition~\ref{prop:T3-local-iterate} first with \(t=\delta\log3\) and then with
\(t=\delta\log2\); the leading constants cancel in the ratio.
\end{proof}

The striking contact \(2^k\) in each individual error does \emph{not} give approximation
order \(2^k\) for the ratio.  The first relative correction is always quadratic in
\(\delta\).  Increasing \(k\) changes the target from \(\thetaAE^2\) to
\(\thetaAE^{2^k}\), and eventually enlarges the correction coefficient; it does not turn a
continued-fraction error into a super-Dirichlet approximation.

\subsection{Height accounting}

For convergents \(p/q\), the logarithmic heights of \(z_2,z_3\) are \(O_{M,A}(q)\).  By
Theorem~\ref{thm:basin-height},
\begin{equation}
 h\bigl(T_3^{\circ k}(z_i)\bigr)=O_{M,A}(2^kq)
 \qquad(i=2,3).
 \label{eq:tower-height}
\end{equation}
The same upper scale holds for the rational ratio \(R_{k;p,q}\).  Continued fractions give
\(|\delta|<1/q_{\mathrm{next}}\), but the trace tower has improved neither the order in
\(\delta\) nor the exponential height scale.  More decisively, for fixed \(M\), a standard
linear-form-in-logarithms estimate applied to \(q\log M-p\log2\) gives an effective lower
bound \(|\delta|\gg_{M}q^{-C_M}\) for some constant \(C_M\)
\cite{BakerWustholz2007}.  Thus the error in~\eqref{eq:tower-approximation} is only of
polynomial scale in \(q\), whereas the value height is exponential in \(q\) and in \(2^k\).
No ordinary rational-separation argument closes that gap.

There is an even more exact obstruction.  Write the initial reduced defect numerators of
\(z_2,z_3\) as \(u_2,u_3\).  Corollary~\ref{cor:no-new-support} says that cubic iteration
creates no common prime outside three.  Hence any exceptional cancellation in
\(R_{k;p,q}\) must already be encoded in the original synchronized binomial gaps.  Iteration
only raises their common valuations to powers.

\begin{target}[What would make the cubic trace useful]
A successful cubic-trace argument must supply a genuinely cross-branch theorem about the
initial reduced gaps of \(M^q/2^p\) and \(A^q/3^p\).  Neither the real superattraction nor
the iteration itself can create that theorem: the exact valuation law proves that all
nonstructural shared support is inherited from stage zero.
\end{target}

This negative conclusion is useful because it is exact.  It rules out an entire class of
``iterate until the rational number is too close to one'' arguments without relying on a
loose height estimate.

\section{Odd prime traces and structural-prime normalization}
\label{sec:prime-traces}

The cubic trace is exceptional, but the prime traces admit another exact theorem which
prepares the fifth-trace bridge.

Let \(p=2h+1\) be an odd prime.  Keep the notation \(S_p,P_p,g_p\) from
Sections~\ref{sec:rational-arithmetic}--\ref{sec:T3}.

\begin{theorem}[Prime-trace congruence]
\label{thm:prime-trace-congruence}
In \(\mathbb F_p[a,b]\),
\begin{equation}
 \boxed{P_p(a,b)\equiv(a-b)^{p-3}\pmod p.}
 \label{eq:Pp-mod-p}
\end{equation}
For coprime positive integers \(a,b\),
\begin{equation}
 g_p(a,b)=
 \begin{cases}
 p,&p\mid a-b,\\
 1,&p\nmid a-b.
 \end{cases}
 \label{eq:gp-prime}
\end{equation}
If \(p\ge5\) and \(p\mid a-b\), then
\begin{equation}
 v_p(P_p(a,b))=1.
 \label{eq:Pp-exact-valuation}
\end{equation}
\end{theorem}

\begin{proof}
Freshman's dream gives
\[
 S_p(a,b)=\frac{a^p-b^p}{a-b}\equiv(a-b)^{p-1}\pmod p.
\]
Reducing~\eqref{eq:defect-factorization} modulo \(p\) and canceling the polynomial
\((a-b)^2\) proves~\eqref{eq:Pp-mod-p}.  The same congruence for \(S_p\) proves
\eqref{eq:gp-prime}.

Suppose now \(p\ge5\) and \(a\equiv b\pmod p\).  Coprimality makes \(b\) a \(p\)-adic
unit.  At the diagonal,
\begin{align*}
 P_p(b,b)
 &=b^{p-3}\sum_{j=1}^{h}j^2
 =b^{p-3}\frac{h(h+1)(2h+1)}6\\
 &=p\,b^{p-3}\frac{p^2-1}{24},
\end{align*}
which has exact \(p\)-adic valuation one.  By~\eqref{eq:Pp-mod-p}, the first derivative of
\(P_p(a,b)\) with respect to \(a\) vanishes on the diagonal modulo \(p\) (the exponent
\(p-3\) is at least two).  Hence replacing \(a=b\) by any \(a\equiv b\pmod p\) changes
\(P_p\) only modulo \(p^2\).  Its valuation remains exactly one.
\end{proof}

\begin{corollary}[Normalized new prime-trace factor]
\label{cor:normalized-prime-factor}
For \(p\ge5\), define
\begin{equation}
 C_p(a,b)=\frac{P_p(a,b)}{p^{\mathbf1_{p\mid a-b}}}.
 \label{eq:Cp}
\end{equation}
Then \(C_p(a,b)\in\N\), \(p\nmid C_p(a,b)\), and if
\(T_p(a/b)=a'/b'\) is reduced, its defect satisfies
\begin{equation}
 a'-b'=(a-b)^2C_p(a,b).
 \label{eq:prime-reduced-defect}
\end{equation}
\end{corollary}

\begin{proof}
Theorem~\ref{thm:prime-trace-congruence} shows that the structural reduction factor
\(g_p\) is precisely the one factor of \(p\) carried by \(P_p\).  Substitute into
\eqref{eq:reduced-defect-general}.
\end{proof}

At \(p=3\), \(P_3=1\) and the structural factor is instead removed directly from the square
\((a-b)^2\), producing the exceptional recurrence of Section~\ref{sec:T3}.  At every prime
\(p\ge5\), the structural prime is canceled out of the new factor, leaving the old defect
valuation to double without loss.

\subsection{The first three levels}

The prime-trace hierarchy now has a clean interpretation:
\begin{center}
\small
\begin{tabular}{@{}c c c c@{}}
\toprule
Trace & Map degree & New normalized factor & Arithmetic character\\
\midrule
\(T_3\)&2&1&pure defect squaring; no new support\\
\(T_5\)&4&normalized \(a^2+3ab+b^2\)&one reciprocal quadratic; two channels\\
\(T_7\)&6&normalized reciprocal quartic&extension-field channels already appear\\
\bottomrule
\end{tabular}
\end{center}

The fifth trace is therefore the first place where a positive trace can introduce new defect
primes, and also the last place where the new factor is a single rational reciprocal
quadratic.  That combination makes it the natural next object to audit exactly.

\section{The fifth trace and the two cross-gap channels}
\label{sec:T5}

\subsection{The normalized fifth-trace factor}

Set
\begin{equation}
 F(a,b)=a^2+3ab+b^2,
 \qquad
 C(a,b)=\frac{F(a,b)}{5^{\mathbf1_{5\mid a-b}}}.
 \label{eq:F-C}
\end{equation}
By Corollary~\ref{cor:normalized-prime-factor}, \(C(a,b)\) is an integer not divisible by
five.  The exact fifth-trace defect is
\begin{equation}
 T_5(a/b)-1
 =\frac{(a-b)^2F(a,b)}{5a^2b^2},
 \label{eq:T5-defect-unreduced}
\end{equation}
and after fraction reduction,
\begin{equation}
 a'-b'=(a-b)^2C(a,b).
 \label{eq:T5-defect-reduced}
\end{equation}

Thus \(C(a,b)\) is precisely the genuinely new prime-support contribution of the first
noncritical trace.

\subsection{A reciprocal-quadratic identity}

The key algebraic identity is more general than the coefficient three.

\begin{lemma}[Reciprocal quadratic bridge]
\label{lem:quadratic-bridge}
For
\[
 F_\lambda(a,b)=a^2+\lambda ab+b^2
\]
and arbitrary commutative-ring elements \(a,b,c,d,\lambda\),
\begin{equation}
 \boxed{
 cdF_\lambda(a,b)-abF_\lambda(c,d)
 =(ad-bc)(ac-bd).}
 \label{eq:quadratic-bridge}
\end{equation}
\end{lemma}

\begin{proof}
Put formally \(r=a/b\), \(s=c/d\).  The identity is the homogeneous clearing of
\[
 s(r^2+\lambda r+1)-r(s^2+\lambda s+1)
 =(r-s)(rs-1).
\]
A direct expansion proves it over an arbitrary commutative ring, without requiring division.
\end{proof}

The coefficient \(\lambda\) disappears.  Geometrically, two roots of the same reciprocal
quadratic are either equal or reciprocal.  Arithmetically, the two possibilities become the
two factors on the right of~\eqref{eq:quadratic-bridge}.

\subsection{Exact channel separation}

\begin{theorem}[Fifth-trace common-content decomposition]
\label{thm:T5-channel-decomposition}
Let \(a,b,c,d\in\N\) with \(\gcd(a,b)=\gcd(c,d)=1\).  Put
\begin{align}
 C_1&=C(a,b),& C_2&=C(c,d),&G&=\gcd(C_1,C_2),\label{eq:channel-G}\\
 D_-&=ad-bc,&D_+&=ac-bd,\label{eq:channel-D}
\end{align}
and
\begin{equation}
 G_-:=\gcd(G,D_-),
 \qquad
 G_+:=\gcd(G,D_+).
 \label{eq:Gpm}
\end{equation}
Then
\begin{equation}
 \boxed{G=G_-G_+,\qquad\gcd(G_-,G_+)=1.}
 \label{eq:channel-product}
\end{equation}
Consequently every prime power common to the two normalized fifth-trace factors belongs, with
its full valuation, to exactly one of the two channels
\[
 D_-=ad-bc\quad\text{or}\quad D_+=ac-bd.
\]
\end{theorem}

\begin{proof}
Let \(\ell^e\mid G\).  Because \(F(a,b)\) is coprime to \(ab\) and \(F(c,d)\) is
coprime to \(cd\), all four coordinates are nonzero modulo \(\ell\).  Identity
\eqref{eq:quadratic-bridge} with \(\lambda=3\) gives
\[
 e\le v_\ell(D_-)+v_\ell(D_+).
\]
If \(\ell\) divided both channels, then in \(\mathbb F_\ell\) the ratios
\(r=a/b\), \(s=c/d\) would satisfy \(r=s\) and \(rs=1\), hence \(r^2=1\).  The equation
\(r^2+3r+1=0\) excludes \(r=-1\), while \(r=1\) forces \(\ell=5\).  But the normalized
factors \(C_1,C_2\) are not divisible by five.  Therefore \(\ell\) divides exactly one of
\(D_-,D_+\), and its full exponent \(e\) lies in that factor.  Taking all primes proves
\eqref{eq:channel-product}.
\end{proof}

Before structural normalization one obtains a compact lcm form.

\begin{corollary}[Unnormalized lcm bound]
\label{cor:T5-lcm}
With \(F(a,b)=a^2+3ab+b^2\),
\begin{equation}
 \boxed{
 \gcd(F(a,b),F(c,d))
 \mid\lcm\bigl(|ad-bc|,|ac-bd|\bigr).}
 \label{eq:T5-lcm}
\end{equation}
Here \(\lcm(0,n)=|n|\), including \(\lcm(0,0)=0\).
\end{corollary}

\begin{proof}
The proof above applies prime by prime except possibly at five.  If five divides
\(F(a,b)\), then \(a\equiv b\pmod5\), and
\[
 F(a,b)=5\left(b^2+5b\frac{a-b}{5}+5\left(\frac{a-b}{5}\right)^2\right).
\]
It automatically has exact five-adic valuation one.  Thus even when both channels are divisible
by five, the lcm contains the full common five-part.
\end{proof}

The product bound follows immediately from Lemma~\ref{lem:quadratic-bridge}; the lcm bound
and the coprime normalized decomposition are the sharper content.  They say that there is no
third, hidden fifth-trace cancellation channel.

\subsection{Specialization to the Alaoglu--Erd\H{o}s germ}

For any sufficiently large convergent, so that \(p,q>0\), write the synchronized rational
near-units in lowest terms as
\begin{equation}
 \begin{aligned}
 z_2&=\frac ab=\frac{M^q/g_2}{2^p/g_2},\\
 z_3&=\frac cd=\frac{A^q/g_3}{3^p/g_3}.
 \end{aligned}
 \label{eq:reduced-germ}
\end{equation}
where
\begin{equation}
 g_2=\gcd(M^q,2^p),
 \qquad
 g_3=\gcd(A^q,3^p).
 \label{eq:g2g3}
\end{equation}
Then the two channels of Theorem~\ref{thm:T5-channel-decomposition} are exactly
\begin{align}
 D_-&=\frac{3^pM^q-2^pA^q}{g_2g_3},
 \label{eq:AE-direct-gap}\\
 D_+&=\frac{(MA)^q-6^p}{g_2g_3}.
 \label{eq:AE-reciprocal-gap}
\end{align}
Using \(\delta=q\beta-p\), the cleared gaps have the analytic forms
\begin{align}
 3^pM^q-2^pA^q
 &=6^p(2^\delta-3^\delta)
 =-6^p\log(3/2)\,\delta+O(6^p\delta^2),
 \label{eq:direct-asymptotic}\\
 (MA)^q-6^p
 &=6^p(6^\delta-1)
 =6^p\log6\,\delta+O(6^p\delta^2).
 \label{eq:reciprocal-asymptotic}
\end{align}

\begin{theorem}[Two-channel reduction for a counterexample]
\label{thm:AE-two-channel}
Under the hypothetical counterexample~\eqref{eq:hypothesis}, every common prime power of the
new normalized fifth-trace factors
\[
 C(M^q/g_2,2^p/g_2),
 \qquad
 C(A^q/g_3,3^p/g_3)
\]
is carried, with its full valuation, by exactly one of the two explicit integers
\eqref{eq:AE-direct-gap} and~\eqref{eq:AE-reciprocal-gap}.
\end{theorem}

\begin{proof}
Apply Theorem~\ref{thm:T5-channel-decomposition} to the reduced fractions
\eqref{eq:reduced-germ}; equations~\eqref{eq:AE-direct-gap}--\eqref{eq:AE-reciprocal-gap}
are direct substitutions.
\end{proof}

This is the central new reduction.  The first trace capable of introducing new defect primes
reduces all shared new support to two synchronized binomial gaps.  Those gaps are still large in absolute value: the same linear-form estimate used above
makes \(|\delta|\) at worst polynomially small in \(q\), while \(6^p\) is exponential.
Thus the theorem is not by itself a contradiction.  It does, however,
replace an amorphous cross-gcd problem by two named and structurally different channels:
\begin{itemize}
\item \(D_-\) compares the two branches in the same direction, \(z_2-z_3\);
\item \(D_+\) compares a branch with the reciprocal of the other, \(z_2z_3-1\).
\end{itemize}
Any Kummer, support, or \(S\)-unit theorem applied to the fifth trace can now be tested against
these exact channels rather than against the full trace factors.

\subsection{Small exact examples}

Table~\ref{tab:T5-channels} illustrates all three possibilities: purely reciprocal, purely
direct, and a split common factor.  Every entry is an exact integer calculation.

\begin{table}[htbp]
\centering
\small
\begin{tabular}{@{}c c r r r@{}}
\toprule
\(a/b\) & \(c/d\) & \(\gcd(C_1,C_2)\) & \(D_-=ad-bc\) & \(D_+=ac-bd\)\\
\midrule
\(37/38\)&\(41/42\)&\(79\)&\(-4\)&\(-79\)\\
\(38/37\)&\(41/42\)&\(79\)&\(79\)&\(4\)\\
\(35/37\)&\(64/65\)&\(11\cdot31\)&\(-3\cdot31\)&\(-3\cdot5\cdot11\)\\
\(18/19\)&\(101/105\)&\(29\cdot59\)&\(-29\)&\(-3\cdot59\)\\
\bottomrule
\end{tabular}
\caption{Exact fifth-trace channel allocation.  The normalized factors equal the displayed
quadratics in these examples because none is divisible by five.}
\label{tab:T5-channels}
\end{table}

For example,
\[
 F(35,37)=6479=11\cdot19\cdot31,
 \qquad
 F(64,65)=20801=11\cdot31\cdot61.
\]
The common prime \(31\) is direct and the common prime \(11\) is reciprocal.  Theorem
\ref{thm:T5-channel-decomposition} predicts this allocation with full valuations, not merely
at the radical level.

\subsection{What remains missing}

Theorem~\ref{thm:AE-two-channel} is an upper-structure theorem: it identifies where common
trace content is allowed to go.  It does not force the dyadic and triadic trace factors to
have a large common divisor in the first place.  A complete proof by this route needs a second
ingredient of one of the following types.

\begin{target}[Forced shared trace content]
Prove, along a suitable infinite sequence of approximants \((p,q)\), a lower bound for the
common part of the two normalized fifth-trace factors that cannot be accommodated by the
direct and reciprocal gap channels.
\end{target}

\begin{target}[Channel-specific incompatibility]
Prove that primes generated by one of the two normalized reciprocal quadratics have a local
or Kummer signature incompatible with both
\(3^pM^q-2^pA^q\) and \((MA)^q-6^p\), except for a quantitatively insufficient exceptional
set.
\end{target}

\begin{target}[Higher trace with controlled resultants]
For \(p\ge7\), common roots of \(P_p(X,1)\) have more than the equal/reciprocal alternatives.
Compute the corresponding bivariate resultants and determine whether their additional
channels carry stronger order or splitting restrictions than the fifth-trace pair.
\end{target}

The first target is the hardest: real synchronization alone does not transparently force
modular synchronization.  The second and third are more local and may be accessible to exact
finite-field, resultant, or Chebotarev analysis.  The fifth trace is valuable because it
provides the smallest complete test case for that program.

\section{Signed trace filters: higher contact and its exact cost}
\label{sec:signed-filters}

Positive traces cannot beat the contact/degree barrier.  Signed products can cancel initial
Bernoulli moments, so it is natural to test them separately.

Let \(d_1,\ldots,d_m\ge3\) be distinct odd integers and
\(c_1,\ldots,c_m\in\Z\).  Define the rational function
\begin{equation}
 \mathcal F(z)=\prod_{i=1}^{m}T_{d_i}(z)^{c_i}.
 \label{eq:signed-filter}
\end{equation}
By Theorem~\ref{thm:bernoulli-trace},
\begin{equation}
 \log\mathcal F(e^t)=
 \sum_{r\ge1}\frac{B_{2r}}{2r(2r)!}
 \left(\sum_{i=1}^{m}c_i(d_i^{2r}-1)\right)t^{2r}.
 \label{eq:signed-moments}
\end{equation}

\begin{proposition}[Moment conditions and support floor]
\label{prop:signed-support}
The filter has contact order at least \(2R\) at one exactly when
\begin{equation}
 \sum_{i=1}^{m}c_i(d_i^{2r}-1)=0
 \qquad(1\le r<R).
 \label{eq:moment-conditions}
\end{equation}
If the coefficients are not all zero, these conditions require \(m\ge R\).
\end{proposition}

\begin{proof}
The first assertion is~\eqref{eq:signed-moments}.  If \(m<R\), take the first \(m\)
conditions.  With \(x_i=d_i^2\), their coefficient matrix is
\((x_i^r-1)_{1\le r,i\le m}\).  Factoring \(x_i-1\) from column \(i\) leaves an ordinary
Vandermonde-type evaluation matrix for polynomials of degrees zero through \(m-1\), hence an
invertible matrix.  All \(c_i\) would vanish.
\end{proof}

The cyclotomic factorization makes cancellation among trace factors exact.  If
\begin{equation}
 e_n=\sum_{i:n\mid d_i}c_i,
 \label{eq:cyclotomic-exponents}
\end{equation}
then, up to a rational scalar,
\begin{equation}
 \mathcal F(z)=\prod_{n>1}\mathscr C_n(z)^{e_n}.
 \label{eq:signed-cyclotomic}
\end{equation}
Thus the fully reduced cyclotomic degree budget is measured by
\begin{equation}
 \mathcal D(\mathcal F)=\sum_{n>1}|e_n|\varphi(n).
 \label{eq:cyclotomic-degree-budget}
\end{equation}
The exact rational contact theorem in the supplied report already implies that fixed rational
filters cannot create unbounded order for free; formulas
\eqref{eq:moment-conditions}--\eqref{eq:cyclotomic-degree-budget} identify the cost inside this
specific trace family.

Two low-order examples are instructive:
\begin{align}
 \log\frac{T_3(e^t)^3}{T_5(e^t)}
 &=\frac{2}{15}t^4-\frac{2}{27}t^6+O(t^8),
 &\mathcal D&=10,
 \label{eq:filter-order4}\\
 \log\frac{T_3(e^t)^8T_9(e^t)}{T_7(e^t)^3}
 &=\frac{64}{63}t^6+O(t^8),
 &\mathcal D&=42.
 \label{eq:filter-order6}
\end{align}
The contact rises from two to four and six, but the exact cyclotomic degree budget rises much
faster.  Signed filters therefore reproduce, in multiplicative trace language, the same
finite-order height obstruction found by rational Pad\'e and Richardson constructions in the
unified report.  Their value here is diagnostic: the Bernoulli moment equations and
cyclotomic degree reveal exactly where the cost enters.

\begin{warning}
No universal lower bound for the \emph{reduced value height} of every signed filter is claimed
here; special evaluations can have accidental arithmetic cancellation.  The exact statements
are the moment conditions, the support floor, and the cyclotomic factorization.  Any proposed
signed filter must still exhibit a branch-specific cancellation theorem rather than relying
on formal contact alone.
\end{warning}

\section{Synthesis: what has been gained and what has not}
\label{sec:synthesis}

\subsection{The new exact statements}

The trace investigation produces the following chain of results, none of which is a restatement
of the determinant or secant arguments in the supplied report.

\begin{enumerate}
\item \textbf{Centered cyclotomic resolution.}
The positive trace \(T_d\) factors into centered cyclotomic factors and has an exact
Bernoulli--Jordan expansion.

\item \textbf{Complete rational reduction.}
For a reduced fraction, cancellation in one trace is bounded by the fixed integer \(d\), and
the defect has the positive sum-of-squares factor \(P_d\).

\item \textbf{Unique criticality.}
The local contact/degree ratio of every positive trace composition is at most one, with equality
only for pure \(T_3\) iteration.

\item \textbf{Closed cubic dynamics.}
The \(T_3\) defect numerator has a complete prime-valuation formula; no new nonstructural prime
can enter, and rational height grows on exactly the contact scale.

\item \textbf{Synchronized tower diagnosis.}
A counterexample would generate rational approximants to \(\thetaAE^{2^k}\), but their
relative error remains quadratic in the original logarithmic gap, while height grows like
\(2^kq\).

\item \textbf{Prime-trace structural normalization.}
For every odd prime \(p\ge5\), the unique structural factor \(p\) can be removed from
\(P_p\), leaving a \(p\)-free new defect factor.

\item \textbf{Two-channel fifth-trace theorem.}
Every prime power shared by the two new normalized fifth-trace factors belongs uniquely to the
direct cross-gap or the reciprocal cross-gap.  In the counterexample specialization these are
\(3^pM^q-2^pA^q\) and \((MA)^q-6^p\).
\end{enumerate}

\subsection{Why this does not yet prove the conjecture}

The cubic map is too perfect in the wrong way: its real contraction and height growth are
exactly balanced, and its valuation law proves that it cannot create missing shared support.
The fifth map creates new support, but no present theorem forces the two branches to share a
large part of that support.  The channel decomposition tells us where a forced common factor
would have to land, not that one exists.

This is not merely the generic four-exponentials wall stated in new notation.  It is a finer
arithmetic diagnosis inside a concrete rational dynamical system.  A future argument can now
fail or succeed at one of three explicit interfaces:
\begin{enumerate}
\item force common normalized trace support;
\item forbid that support in both cross-gap channels;
\item move to a higher prime trace and exploit a genuinely stronger resultant channel.
\end{enumerate}
There is no longer any ambiguity about what repeated cubic contraction or a single fifth
trace can accomplish by themselves.

\subsection{Prioritized continuation}

\begin{target}[Priority A: formalize the exact trace core]
The elementary core consists of
Theorems~\ref{thm:exact-denominator} and \ref{thm:defect-factorization}, the cubic results
\ref{thm:T3-recurrence} and \ref{thm:T3-valuations}, the prime-trace result
\ref{thm:prime-trace-congruence}, and the fifth-trace result
\ref{thm:T5-channel-decomposition}.  These use only finite sums, gcds, polynomial identities,
and elementary valuations.  They are suitable for immediate Lean
formalization, independent of any unformalized transcendence theorem.
\end{target}

\begin{target}[Priority B: exact gcd reconnaissance on the two AE channels]
For candidate-like integer pairs \((M,A)\) chosen near \(A=M^{\thetaAE}\), compute the
normalized fifth-trace factors and their direct/reciprocal allocation over many convergents.
The purpose is not probabilistic evidence for the conjecture; it is to identify whether one
channel has a stable valuation or splitting signature worth proving.
\end{target}

\begin{target}[Priority C: the seventh-trace resultant]
Factor \(P_7(X,1)=X^4+3X^3+6X^2+3X+1\) over its natural quadratic field, compute the
bivariate common-root resultant, and classify the finite-field channels by Frobenius type.
The fifth trace has only equal/reciprocal channels; the seventh is the first test of whether
extension-field channels impose stronger prime-order conditions.
\end{target}

\begin{target}[Priority D: channel-sensitive \(S\)-unit bounds]
The direct and reciprocal gaps have different multiplicative shapes.  Seek gcd or primitive-
divisor estimates which distinguish
\(3^pM^q-2^pA^q\) from \((MA)^q-6^p\), rather than applying a symmetric generic bound to
their product.  The channel theorem shows that such an asymmetric estimate would be used
without loss.
\end{target}

\section{Lean-oriented dependency map}
\label{sec:lean}

The elementary core can be formalized in small layers.  Suggested theorem names are included
only as an implementation blueprint.

\begin{longtable}{@{}p{0.28\textwidth}p{0.33\textwidth}p{0.31\textwidth}@{}}
\caption{Suggested formalization order.}\label{tab:lean-map}\\
\toprule
Module & Principal statements & Dependencies\\
\midrule
\endfirsthead
\toprule
Module & Principal statements & Dependencies\\
\midrule
\endhead
\texttt{CenteredTrace.Basic}
& \texttt{centeredTrace}, reciprocity, geometric-sum form, positivity
& finite sums, ordered fields, AM--GM for a finite positive list\\
\addlinespace
\texttt{CenteredTrace.Rational}
& exact formula \eqref{eq:rational-Td}, \texttt{gcd\_Sd\_mul}, reduction factor
& gcd and modular arithmetic for naturals\\
\addlinespace
\texttt{CenteredTrace.Defect}
& sum-of-squares \(P_d\), factorization \eqref{eq:defect-factorization}, positivity
& polynomial identities, finite reindexing\\
\addlinespace
\texttt{CenteredTrace.Cubic}
& reduced recurrence, \(g_k=3\iff3\mid u_k\), closed valuation law
& previous module, elementary \(p\)-adic valuations\\
\addlinespace
\texttt{CenteredTrace.Height}
& bounds \eqref{eq:height-growth}, iterate bounds, canonical-height limit
& integer inequalities, geometric sequences, Cauchy criterion\\
\addlinespace
\texttt{CenteredTrace.Prime}
& \(P_p\equiv(a-b)^{p-3}\), exact structural normalization
& finite-field Frobenius, sum of squares formula\\
\addlinespace
\texttt{CenteredTrace.Fifth}
& bridge identity, lcm bound, normalized coprime channel decomposition
& gcd valuations and finite-field division by units\\
\addlinespace
\texttt{CenteredTrace.Cyclotomic}
& factorization \eqref{eq:trace-cyclotomic-factorization}
& cyclotomic polynomial API; can be postponed\\
\bottomrule
\end{longtable}

The critical proof engineering choice is to formalize the homogeneous integer identities
before introducing rational numbers.  For example, the fifth-trace bridge should first be
entered literally as
\[
 cd(a^2+3ab+b^2)-ab(c^2+3cd+d^2)=(ad-bc)(ac-bd),
\]
then specialized to fractions only after all divisibility consequences are available.
Likewise, the cubic recurrence is cleaner as a theorem about coprime coordinate pairs than as
a theorem about normalized rationals.

\begin{statusbox}{No Lean acceptance is claimed for this addendum.  The map is intentionally
split into elementary lemmas whose statements avoid analytic branches, real logarithms, and
external transcendence inputs.}
\end{statusbox}

\section{Conclusion}

The centered trace family began as an attempt to exploit the synchronized rational near-units
produced by a hypothetical Alaoglu--Erd\H{o}s counterexample.  It succeeded in producing an
exact superattractor, but the arithmetic audit is as important as the contraction:
\begin{itemize}
\item \(T_3\) squares the defect exactly and doubles its valuations, yet its rational height
has exactly the same dynamical degree;
\item iteration cannot manufacture a missing common prime between the dyadic and triadic
branches;
\item \(T_5\) is the first level that introduces new defect support, and that support has a
complete two-channel common-divisor law;
\item all shared normalized fifth-trace content is carried by two explicit synchronized
binomial gaps.
\end{itemize}

The conjecture remains unclosed, but the route has not merely been tried and abandoned.  It
has been reduced to an exact arithmetic interface.  Any future trace-based proof must obtain
new information about shared support or about the two channel gaps; no amount of unaudited
iteration, higher local contact, or hidden single-branch denominator cancellation can replace
that information.

\appendix

\section{Exact verification script}
\label{app:code}

The following standard-library Python program checks the finite polynomial identities,
reduced cubic recurrence, valuation formulas on a test box, and the fifth-trace channel
decomposition.  It uses no floating-point arithmetic for theorem checks.  The decimal orbit
displays are derived only after the exact fractions have been formed.

\begin{lstlisting}[style=pythonstyle,caption={Exact centered-trace verifier.}]
from fractions import Fraction
from itertools import product
from math import gcd, lcm, log10


def S(d: int, a: int, b: int) -> int:
    return sum(a ** (d - 1 - j) * b ** j for j in range(d))


def P(d: int, a: int, b: int) -> int:
    h = (d - 1) // 2
    total = 0
    for j in range(1, h + 1):
        q = sum(a ** (j - 1 - r) * b ** r for r in range(j))
        total += (a * b) ** (h - j) * q * q
    return total


def trace_step(d: int, a: int, b: int) -> tuple[int, int, int]:
    h = (d - 1) // 2
    num = S(d, a, b)
    den = d * (a * b) ** h
    g = gcd(num, den)
    return num // g, den // g, g


def vp(n: int, p: int) -> int:
    n = abs(n)
    out = 0
    while n and n % p == 0:
        n //= p
        out += 1
    return out


# General defect and exact-denominator identities.
for d in range(3, 20, 2):
    h = (d - 1) // 2
    for a, b in product(range(1, 30), repeat=2):
        if gcd(a, b) != 1:
            continue
        assert gcd(S(d, a, b), a * b) == 1
        assert S(d, a, b) - d * (a * b) ** h == (a - b) ** 2 * P(d, a, b)
        aa, bb, g = trace_step(d, a, b)
        assert g == gcd(S(d, a, b), d)
        assert aa - bb == (a - b) ** 2 * P(d, a, b) // g


# Cubic closed form and valuation transport.
for a, b in product(range(1, 80), repeat=2):
    if gcd(a, b) != 1 or a == b:
        continue
    u0 = a - b
    aa, bb = a, b
    for k in range(1, 7):
        aa, bb, g = trace_step(3, aa, bb)
        uk = aa - bb
        if u0 % 3:
            assert abs(uk) == abs(u0) ** (2 ** k)
        else:
            assert abs(uk) == 3 * (abs(u0) // 3) ** (2 ** k)
        for ell in (2, 5, 7, 11, 13):
            assert vp(uk, ell) == (2 ** k) * vp(u0, ell)


# Fifth-trace bridge and lcm/channel theorem.
def F(a: int, b: int) -> int:
    return a * a + 3 * a * b + b * b


def C5(a: int, b: int) -> int:
    return F(a, b) // (5 if (a - b) % 5 == 0 else 1)


for a, b, c, d in product(range(1, 45), repeat=4):
    if gcd(a, b) != 1 or gcd(c, d) != 1:
        continue
    dm = a * d - b * c
    dp = a * c - b * d
    assert c * d * F(a, b) - a * b * F(c, d) == dm * dp

    raw = gcd(F(a, b), F(c, d))
    assert lcm(abs(dm), abs(dp)) % raw == 0

    common = gcd(C5(a, b), C5(c, d))
    gm = gcd(common, abs(dm))
    gp = gcd(common, abs(dp))
    assert gcd(gm, gp) == 1
    assert gm * gp == common

print("all exact checks passed")
\end{lstlisting}

\section{Proof details for the coefficient formula}
\label{app:coefficients}

For completeness, expand
\[
 P_{2h+1}(a,b)=\sum_{j=1}^{h}(ab)^{h-j}
 \left(\sum_{s=0}^{j-1}a^{j-1-s}b^s\right)^2.
\]
Fix the exponent \(r\) of \(b\), with \(0\le r\le2h-2\).  In the \(j\)-th summand, the
base factor contributes \(h-j\), and the square contributes the number of ordered pairs
\((s,t)\in\{0,\ldots,j-1\}^2\) with
\(s+t=r-h+j\).  Summing this triangular count over \(j\) gives
\[
 \binom{\min(r,2h-2-r)+2}{2},
\]
which proves~\eqref{eq:Pd-coefficients}.  The coefficient sequence is therefore
\[
 1,3,6,10,\ldots,\binom{h+1}{2},\ldots,10,6,3,1.
\]
The positivity visible in the sum-of-squares form is thus also coefficientwise.

\section{Status matrix}
\label{app:status}

\begin{longtable}{@{}>{\raggedright\arraybackslash}p{0.55\textwidth}>{\raggedright\arraybackslash}p{0.16\textwidth}>{\raggedright\arraybackslash}p{0.21\textwidth}@{}}
\caption{Logical status of the principal claims.}\label{tab:status}\\
\toprule
Statement & Status & Location\\
\midrule
\endfirsthead
\toprule
Statement & Status & Location\\
\midrule
\endhead
Synchronized rational near-unit pair from a hypothetical counterexample
&\statusproved&Proposition~\ref{prop:germ}\\
Centered trace positivity, reciprocity, and hyperbolic form
&\statusproved&Proposition~\ref{prop:basic-trace}\\
Bernoulli expansion and centered cyclotomic/Jordan factorization
&\statusproved&Theorems \ref{thm:bernoulli-trace} and \ref{thm:cyclotomic-resolution}\\
Exact rational denominator and positive defect factorization
&\statusproved&Theorems \ref{thm:exact-denominator} and \ref{thm:defect-factorization}\\
Unique contact/degree criticality of pure cubic compositions
&\statusproved&Theorem~\ref{thm:criticality}\\
Closed cubic recurrence and all-prime valuation law
&\statusproved&Theorems \ref{thm:T3-recurrence} and \ref{thm:T3-valuations}\\
Real convergence, height escape, and local iterate asymptotic
&\statusproved&Theorem \ref{thm:basin-height}; Proposition \ref{prop:T3-local-iterate}\\
Rational approximants to \(\thetaAE^{2^k}\) with quadratic relative error
&\statusproved&Theorem~\ref{thm:tower-approximants}\\
Odd-prime structural normalization
&\statusproved&Theorem \ref{thm:prime-trace-congruence}; Corollary \ref{cor:normalized-prime-factor}\\
Fifth-trace lcm and coprime two-channel decomposition
&\statusproved&Theorem \ref{thm:T5-channel-decomposition}; Corollary \ref{cor:T5-lcm}\\
Exact finite test tables and regression checks
&\statuscomputed&Table~\ref{tab:T3-orbits}; Appendix~\ref{app:code}\\
A lower bound forcing shared fifth-trace support
&\statusopen&Section~\ref{sec:T5}\\
A channel-specific Kummer or \(S\)-unit contradiction
&\statusopen&Section~\ref{sec:synthesis}\\
Lean formalization of the new elementary core
&\statusopen&Section~\ref{sec:lean}\\
\bottomrule
\end{longtable}

\begin{thebibliography}{99}

\bibitem{UnifiedReport2026}
The ProveIt project,
\newblock \emph{The Alaoglu--Erd\H{o}s Integer-Exponent Conjecture: Machine-Checked Partial
Results, Exact Conditional Structure, Determinant Methods, and the Remaining Barrier},
\newblock unified working report, August 22, 2026.

\bibitem{AE1944}
L.~Alaoglu and P.~Erd\H{o}s,
\newblock \emph{On highly composite and similar numbers},
\newblock Transactions of the American Mathematical Society \textbf{56} (1944), 448--469.

\bibitem{MetropolisRota1983}
N.~Metropolis and G.-C.~Rota,
\newblock \emph{Witt vectors and the algebra of necklaces},
\newblock Advances in Mathematics \textbf{50} (1983), 95--125.

\bibitem{Silverman2007}
J.~H.~Silverman,
\newblock \emph{The Arithmetic of Dynamical Systems},
\newblock Graduate Texts in Mathematics 241, Springer, 2007.

\bibitem{BakerWustholz2007}
A.~Baker and G.~W\"ustholz,
\newblock \emph{Logarithmic Forms and Diophantine Geometry},
\newblock New Mathematical Monographs 9, Cambridge University Press, 2007.

\bibitem{BCZ2003}
Y.~Bugeaud, P.~Corvaja, and U.~Zannier,
\newblock \emph{An upper bound for the G.C.D. of \(a^n-1\) and \(b^n-1\)},
\newblock Mathematische Zeitschrift \textbf{243} (2003), 79--84.

\end{thebibliography}

\end{document}
